\documentclass[12pt, a4paper]{article}
\usepackage[utf8]{inputenc}
\usepackage[T1]{fontenc}
\usepackage[portuguese]{babel}
\usepackage{geometry}
\geometry{a4paper, margin=2.5cm}
\begin{document}

% TÍTULO
\title{Resumo Executivo: Identificação de Locais Ótimos e Tecnologias para Sistemas Solares Comunitários em Angola}
\author{Rocélio Da Silva et al.}
\date{\today}
\maketitle

\section*{Resumo}
Este resumo executivo sintetiza os principais pontos do estudo completo sobre a identificação de locais ótimos e tecnologias para implementação de sistemas solares comunitários em Angola, com foco no framework GEESP-Angola e estudo de caso na província da Huíla.

\textbf{Contexto:} Angola apresenta baixa taxa de eletrificação rural (<10\%), criando barreiras ao desenvolvimento social e econômico. O acesso à energia é fundamental para saúde, educação e produtividade.

\textbf{Objetivo:} Desenvolver e aplicar um framework de análise multicritério (MCDA-GIS) para selecionar sítios e tecnologias solares apropriadas, integrando critérios técnicos, socioeconômicos e de infraestrutura.

\textbf{Metodologia:} Foram integradas seis camadas de dados (irradiação solar, declividade, densidade populacional, distância à rede, NDVI, luminosidade noturna) em ambiente SIG, com pesos definidos via Analytic Hierarchy Process (AHP). O framework recomenda tecnologias (mini-redes, off-grid, rastreadores, híbridos) conforme perfil de carga, LCOE e contexto local. A robustez foi testada por análise de sensibilidade (±20\% nos pesos).

\textbf{Resultados:} Aplicação na Huíla identificou 45 comunidades e 3 zonas prioritárias (Cacula, Humpata, Quilengues). O índice de aptidão máximo foi 0,83. O LCOE estimado para mini-redes é USD 0,24–0,35/kWh, beneficiando ~95.000 pessoas. Recomendações tecnológicas variam conforme o contexto: mini-redes para vilas maiores, off-grid para comunidades isoladas, rastreadores para áreas com restrição de espaço, híbridos para zonas de transição.

\textbf{Discussão:} O método supera abordagens tradicionais, reduzindo o erro de seleção de sítios em até 75\%. Destaca-se a importância da governança local, capacitação e validação em campo. O framework é replicável para outras províncias e países africanos.

\textbf{Logística e Implementação:} O sucesso depende de planejamento detalhado, parcerias estratégicas, financiamento diversificado e gestão de riscos logísticos. Recomenda-se integração com políticas públicas e capacitação local.

\textbf{Conclusão:} O GEESP-Angola oferece uma abordagem baseada em evidências para universalização da energia solar comunitária, promovendo justiça energética e resiliência climática. Todos os dados, códigos e protocolos estão disponíveis em repositório público para replicação.

\end{document}
